% =====================================================================
% sec/1_intro.tex  --  Introduction
% =====================================================================

\section{Introduction}
\label{sec:intro}

Colorectal cancer is one of the leading causes of cancer-related mortality worldwide.
Early detection and accurate delineation of colorectal polyps during colonoscopy are
essential for reducing mortality. While deep learning has made remarkable progress in
medical image segmentation, generalizing across diverse colonoscopy equipment and
imaging protocols remains a fundamental challenge.

\paragraph{Motivation.}
Recent vision-language pre-trained (VLP) models such as CLIP~\cite{radford2021clip}
and MedCLIP~\cite{wang2022medclip} have demonstrated strong cross-modal alignment
between images and text. However, directly adapting general-domain VLP models to
colorectal polyp segmentation introduces two limitations:
\begin{enumerate}
    \item \textbf{Domain gap}: General VLP visual encoders are not exposed to
          endoscopic image statistics, leading to suboptimal feature representations.
    \item \textbf{Deterministic adaptation}: Standard fine-tuning produces point
          estimates without capturing epistemic uncertainty, which is critical for
          clinical decision support.
\end{enumerate}

\paragraph{Contributions.}
We propose \textbf{ProbVLP}, a probabilistic adaptation framework for colorectal
polyp segmentation that addresses both limitations:
\begin{itemize}
    \item \textbf{EndoViT visual backbone}: We replace the CLIP visual tower with
          EndoViT~\cite{} -- a ViT-B/16 pretrained via masked autoencoding on
          Endo700k -- to obtain endoscopy-specific visual features.
    \item \textbf{LoRA-adapted Probabilistic Vision-Language (PVL) adapters}: We
          inject low-rank adaptation (LoRA) into the bidirectional cross-attention
          layers of the PVL adapters, enabling parameter-efficient modality alignment
          with as few as $r{=}8$ trainable parameters per projection matrix.
    \item \textbf{Boundary-feedback gating}: A lightweight boundary prediction
          branch provides spatial priors that are gated back into subsequent PVL
          adapter forward passes via learnable per-adapter scalar gates.
    \item \textbf{Inverse-variance TTA}: At inference time, we fuse predictions
          from multiple augmented views (flips + scales + MC-dropout samples) using
          inverse-variance weighting, producing calibrated uncertainty maps alongside
          segmentation masks.
\end{itemize}

Trained exclusively on Kvasir-SEG, ProbVLP achieves state-of-the-art cross-domain
generalization to four unseen datasets (ClinicDB, ColonDB, CVC300, BKAI), validating
its applicability in real-world clinical settings.
